\documentclass[11pt]{article}
\usepackage[margin=0.95in]{geometry}
\usepackage{amsmath,amssymb}
\usepackage{graphicx}
\usepackage{booktabs,longtable,array}
\usepackage{xcolor}
\usepackage[hidelinks]{hyperref}
\usepackage{enumitem}
\usepackage{caption}
\setlist{itemsep=2pt,topsep=4pt}
\captionsetup{font=small,labelfont=bf}
\setlength{\parskip}{5pt}
\setlength{\parindent}{0pt}
\renewcommand{\arraystretch}{1.15}
\title{Glossary: Every Label, Code and Abbreviation in the Project}
\author{}
\date{trm-delay-robust-scaling $\cdot$ study reference $\cdot$ 2026-10-08\\[2pt]Covers Phases 0--4. Each entry gives the label, what it means, and where it is used. Gate outcomes are the final recorded ones.}
\begin{document}
\maketitle
\tableofcontents
\clearpage



\section*{0. Read this first: labels that look alike}\addcontentsline{toc}{section}{0. Read this first: labels that look alike}



Several short codes are reused with different meanings in different phases. These are the ones most likely to confuse:



\begin{center}\small\begin{tabular}{>{\raggedright\arraybackslash}p{0.069\linewidth}>{\raggedright\arraybackslash}p{0.563\linewidth}>{\raggedright\arraybackslash}p{0.274\linewidth}}
\toprule
You see & It can mean & How to tell \\
\midrule
\textbf{C1} vs \textbf{C-1} & C1 = a delay \emph{condition} (fixed 1-step delay + queue). C-1 = the \emph{capacity confound} from the Phase 1--2 analysis. & The hyphen. Conditions have none; confounds do. \\
\textbf{A-1} & Phase 3 amendment A-1 (theory check on the headroom probe) \textbf{or} Phase 4 amendment A-1 (Stage A clip norms). & Check the phase. \\
\textbf{A1--A4} & Phase 3 Stage 0.1 headroom \emph{arms} (no hyphen). Also codes A1--A9 in the Phase 3 plan's exploration inventory. & Arms appear in Stage 0 tables; inventory codes only in the plan's \S{}6. \\
\textbf{E1--E5} & Phase 1--2 report: E1--E4 = mechanistic \emph{explanations} (\S{}7), E1--E5 = recommended \emph{experiments} (\S{}10). Phase 3 plan: E-1--E-9 = \emph{engineering items}. Phase 4: \textbf{Stage E} = the Lipschitz analysis. & Context: explanation, experiment, engineering task or stage. \\
\textbf{G1} & Phase 3 \emph{gate} G1 (capacity vs architecture) \textbf{or} inventory item G1 in the Phase 3 plan (out-of-scope: embedded deployment). & Gates are in tables of outcomes; inventory items only in plan \S{}6. \\
\textbf{"788k"} & The parameter count of TRM / MLP \emph{under DDPG} (actor + one critic). Under TD3 (two critics) the same configs total \textasciitilde{}1.18M. & The label names the configuration, not the TD3 total. \\
\textbf{H1} vs \textbf{H1$'$} & H1 = Phase 1's complexity hypothesis. H1$'$ = Phase 4's measurable version of it. & The prime mark. \\
\textbf{H} (alone) & \texttt{H\_cycles}, TRM's number of outer cycles (H = 1, 3, 6). & Appears as "H = 3", "TRM H = 1". \\
\textbf{T} & \textbf{Stage T} in Phase 4 (G1 under TD3) \textbf{or} TRC's number of control steps (T = 50--100). & Stage T is ours; T as a horizon only appears in the TRC comparison. \\
\bottomrule
\end{tabular}\end{center}



\medskip



\section*{1. The project in one line}\addcontentsline{toc}{section}{1. The project in one line}



\begin{center}\small\begin{tabular}{>{\raggedright\arraybackslash}p{0.276\linewidth}>{\raggedright\arraybackslash}p{0.658\linewidth}}
\toprule
Term & Meaning \\
\midrule
\textbf{Research question} & Does weight-tied recursion (TRM), used as the actor-critic trunk of a model-free RL agent, benefit continuous control under action delay --- and if not, why not? \\
\textbf{trm-delay-robust-scaling} & The project's name, the epic under \texttt{epics/}, and the W\&B project. \\
\textbf{Trunk} & The shared body of a network that turns an observation into features. The actor and the critic each have one; a "head" (one linear layer) turns features into an action or a Q-value. The project swaps only the trunk. \\
\textbf{TRM vs MLP} & The central comparison: a weight-tied recursive trunk against a plain feed-forward one. \\
\bottomrule
\end{tabular}\end{center}



\section*{2. Hypotheses and preconditions}\addcontentsline{toc}{section}{2. Hypotheses and preconditions}



\begin{center}\small\begin{tabular}{>{\raggedright\arraybackslash}p{0.161\linewidth}>{\raggedright\arraybackslash}p{0.344\linewidth}>{\raggedright\arraybackslash}p{0.115\linewidth}>{\raggedright\arraybackslash}p{0.258\linewidth}}
\toprule
Label & Statement & Introduced & Outcome \\
\midrule
\textbf{H1} & TRM's advantage over the MLP (and Transformer) is non-decreasing in task complexity; measurable once the horizon is \textasciitilde{}5$\times$ Pendulum's with sparse reward. & Phase 0 ideation, tested Phase 1 & Threshold prediction falsified on MountainCar; complexity scaling as a whole \textbf{untested} (no harder task ever learned by the baseline). \\
\textbf{H1$'$} & Phase 4's operational H1: TRM's deficit, normalised by the random-to-best span, is smaller on a harder task than on Pendulum ($\geq$ 5/6 seed pairs). & Phase 4 plan \S{}7 & \textbf{Not run} (K-3). \\
\textbf{H\_precondition} / \textbf{H\_transfer} & Phase 1--2 report's two readings of the negative: TRM needs preconditions (P1--P3) that the port lacks / recursion transfers and the negative reflects something else. & Phase 1--2 report & Superseded by H3, H\_cap and H\_null in Phase 3. \\
\textbf{P1} & Precondition: \emph{refinement continuity}. Per-step supervision plus a persistent state to refine. The port breaks both halves. & Phase 1--2 report & The best-supported account; Phase 3 G3 tested the "persistent state" half (failed), Phase 4 GA tested the "supervision" half (failed). \\
\textbf{P2} & Precondition: a structured, compositionally correctable answer (a grid can be partly right; a 1-D torque cannot). & Phase 1--2 report & Not manipulable here. \\
\textbf{P3} & Precondition: per-instance adaptive capacity (TRM's per-puzzle embedding). No analogue in replay-based RL. & Phase 1--2 report & Not testable here. \\
\textbf{H3} & Residual-uncertainty hypothesis: recursion earns its cost only where (i) the input leaves uncertainty iteration could resolve \emph{and} (ii) the trunk carries information across steps. & Phase 3 plan \S{}2 & \textbf{Rejected} (G2, G2s fail; G3 fails). \\
\textbf{H3.1} & The deficit is a capacity artefact. & Phase 3 & Rejected (G1: architectural). \\
\textbf{H3.2} & The load-bearing interaction: Carry-Z $-$ MLP $\leq$ 0 with nothing to infer, > 0 where there is. & Phase 3 & Rejected (G2). \\
\textbf{H3.3} & Depth pays only when there is something to infer. & Phase 3 & Not supported: depth inert everywhere. \\
\textbf{H3.4} & Recursion trades return for command regularity (smoother at matched return). & Phase 3 & Supported at n = 12 (0.4 vs 9 reversals/s). \\
\textbf{H3.5} & Recursion degrades more gracefully off-nominal. & Phase 3 & Contradicted (TRM less robust). \\
\textbf{H3.6} & Recursion is real-time feasible at robot rates. & Phase 3 & Quantified: 4.1 ms/decision, 100 Hz yes, 500 Hz no. \\
\textbf{H\_cap} & The deficits track parameter count; recursion is irrelevant. & Phase 3 & Rejected (G1). \\
\textbf{H\_null} & Recursive refinement does not transfer to model-free continuous control under any conditioning. & Phase 3 & What remains, restricted to the setting tested. \\
\emph{(Phase 4, as hypotheses)} & TRM was just untuned (Stage C); restoring the training discipline recovers a benefit (Stage A); the G1 attribution holds under TD3 (Stage T). & Phase 4 & Untuned: no (GC insensitive). Discipline: no (GA fail). TD3: direction holds, but via an optimisation failure (GT pass with caveat). \\
\bottomrule
\end{tabular}\end{center}



\section*{3. Phases and stages}\addcontentsline{toc}{section}{3. Phases and stages}



\subsection*{Phases}\addcontentsline{toc}{subsection}{Phases}



\begin{center}\small\begin{tabular}{>{\raggedright\arraybackslash}p{0.074\linewidth}>{\raggedright\arraybackslash}p{0.153\linewidth}>{\raggedright\arraybackslash}p{0.679\linewidth}}
\toprule
Phase & Dates & What it did \\
\midrule
\textbf{Phase 0} & to 2026-09-08 & Groundwork: pilot, delay wrapper, augmentation fix, TRM port, scoping, literature review, hypothesis ideation. \\
\textbf{Phase 1} & 2026-09-08/09 & Horizon and sparsity: Pendulum-v1 vs MountainCarContinuous-v0, 4 trunks, 48 runs (submitted 3 times: v1, v2, v3; v3 canonical). \\
\textbf{Phase 2} & 2026-09-09/10 & Dimensionality and coupling: BipedalWalker-v3, 6 trunks, 3 conditions, 54 runs; then the post-hoc re-analysis. \\
\textbf{Phase 3} & 2026-09-14 to 09-23 & Confound removal: capacity, carried latent, inferable delay; control-quality battery; 231 jobs. \\
\textbf{Phase 4} & 2026-10-06 to 10-08 & Complexity rung, training discipline, LR sweep, G1 under TD3, Lipschitz analysis, C export. \\
\bottomrule
\end{tabular}\end{center}



\subsection*{Phase 3 stages}\addcontentsline{toc}{subsection}{Phase 3 stages}



\begin{center}\small\begin{tabular}{>{\raggedright\arraybackslash}p{0.222\linewidth}>{\raggedright\arraybackslash}p{0.712\linewidth}}
\toprule
Stage & Question \\
\midrule
\textbf{S0.1} (Stage 0.1) & Headroom probe: is there anything for inference to recover? Arms A1--A4. \\
\textbf{S0.2} (Stage 0.2) & Checkpoint battery over the 102 banked Phase 1--2 checkpoints, plus the latency benchmark. \\
\textbf{Stage 1} & Capacity $\times$ depth: \{MLP, TRM\} $\times$ \{135k, 788k\}, plus H = 1 and H = 6. \\
\textbf{Stage 2} & Uncertainty $\times$ memory: \{MLP, Reset-Z, Carry-Z, carried LSTM\} $\times$ \{C1, C3, C4\}. \\
\textbf{Supplementary (s2x)} & Latent per-episode delay C6 (the G2s test). \\
\textbf{Stage 3} & TD3 cross-check and an unstable plant. \textbf{Not run} (gated on G2). \\
\textbf{Stage 4} & Restore TRM's training discipline. \textbf{Not run} in Phase 3 (gated on G3); became Phase 4 Stage A. \\
\textbf{Stage 5} & Synthesis. \\
\bottomrule
\end{tabular}\end{center}



\subsection*{Phase 4 stages}\addcontentsline{toc}{subsection}{Phase 4 stages}



\begin{center}\small\begin{tabular}{>{\raggedright\arraybackslash}p{0.118\linewidth}>{\raggedright\arraybackslash}p{0.678\linewidth}>{\raggedright\arraybackslash}p{0.110\linewidth}}
\toprule
Stage & Question & Runs \\
\midrule
\textbf{B0} & Complexity pilot: does the MLP learn InvertedDoublePendulum-v5 at 150k / 300k / 500k steps? & 9 + 2 timing \\
\textbf{B0h} & Kill-switch fallback: the same pilot on Hopper-v5. & 9 \\
\textbf{B0d} & IDP diagnostic: C1 at 1M steps; no delay at 500k. & 6 \\
\textbf{B1} & Complexity rung main matrix (4 trunks $\times$ 6 seeds on the harder task). & \textbf{Not run} (K-3) \\
\textbf{C} & LR sweep: TRM 788k, actor $\times$ critic learning-rate multipliers \{0.3, 1, 3\}. & 27 \\
\textbf{A} (+ \textbf{a0} pilot) & Training discipline: truncated vs full-gradient vs full-gradient + deep supervision, at C1 and C4 with previous action. & 4 + 36 \\
\textbf{T} & Phase 3 Stage 1 contrast under TD3. & 48 \\
\textbf{E} (run as \textbf{e2}) & Exploratory: actor Lipschitz constant over the 72 Phase 3 Stage 1 checkpoints. & 6 shards \\
\textbf{D} & Export actors to float32 C headers for an STM32. & local \\
\textbf{smoke} & One GPU job exercising every new code path before any data. & 1 \\
\bottomrule
\end{tabular}\end{center}



\section*{4. Gates (pre-registered decision rules)}\addcontentsline{toc}{section}{4. Gates (pre-registered decision rules)}



\begin{center}\small\begin{tabular}{>{\raggedright\arraybackslash}p{0.051\linewidth}>{\raggedright\arraybackslash}p{0.057\linewidth}>{\raggedright\arraybackslash}p{0.469\linewidth}>{\raggedright\arraybackslash}p{0.300\linewidth}}
\toprule
Gate & Phase & Rule (short) & Outcome \\
\midrule
\textbf{G0} & 3 & Headroom: oracle (A3 $-$ A2) or memory (A2 $-$ A4) headroom > pooled SD in $\geq$ 5/6 seeds. & Pass (memory headroom only) \\
\textbf{G1} & 3 & Attribution: 2$\times$2 capacity $\times$ architecture, "different" = $\geq$ 10/12 pairs share a sign. & Architecture, not capacity \\
\textbf{G2} & 3 & Interaction: [Carry-Z $-$ MLP] improves from C1 to C3 or C4 in $\geq$ 5/6 seeds. & Fail (H3 rejected) \\
\textbf{G2s} & 3 & Latent-delay version of G2 (C6), supplementary. & Fail (rejection stands) \\
\textbf{G3} & 3 & Continuity: Carry-Z beats Reset-Z in $\geq$ 5/6 seeds somewhere. & Fail \\
\textbf{G4} & 3 & Practicality: a non-MLP trunk non-dominated on (return, command rate) and admitting 100 Hz. & Weak pass (literal rule) \\
\textbf{GB0} & 4 & MLP avg-100 $\geq$ 8190 on IDP in $\geq$ 2/3 seeds at some budget. & Fail \\
\textbf{KB0} & 4 & Kill switch: same on Hopper ($\geq$ 3420). & Fail (Hopper dropped) \\
\textbf{GC} & 4 & Any non-default LR cell beats $\times$1/$\times$1 by > pooled SD in 3/3 seeds. & Insensitive (no cell flagged) \\
\textbf{GA} & 4 & A discipline arm beats truncated TRM in $\geq$ 5/6 seeds at C1 or C4. & Fail \\
\textbf{GT} & 4 & TRM 788k $-$ MLP 788k negative in $\geq$ 10/12 pairs under TD3. & Pass, with caveat: TRM 788k does not train under \texttt{td3.yaml} \\
\textbf{H1p} (H1$'$) & 4 & Normalised deficit smaller on the new rung in $\geq$ 5/6 seed pairs. & Not run \\
\textbf{E} & 4 & --- (exploratory, not a gate) & Lipschitz does not explain TRM's noise sensitivity \\
\bottomrule
\end{tabular}\end{center}



\section*{5. Delay conditions and observation design}\addcontentsline{toc}{section}{5. Delay conditions and observation design}



\begin{center}\small\begin{tabular}{>{\raggedright\arraybackslash}p{0.231\linewidth}>{\raggedright\arraybackslash}p{0.703\linewidth}}
\toprule
Label & Meaning \\
\midrule
\textbf{d} & Action delay in steps: the action chosen now reaches the plant d steps later. \\
\textbf{Queue / augmentation} (\texttt{augment\_obs}) & Appending the pending (committed, not-yet-applied) actions to the observation. Restores a sufficient statistic for fixed delays (Katsikopoulos \& Engelbrecht, 2003). The Phase 0 fix. \\
\textbf{Previous action / "pa"} (\texttt{prev\_\allowbreak{}action\_\allowbreak{}obs}) & Also appending the last executed action. Phase 3 Stage 2 (amendment A-6) and Phase 4 Stage A only. \\
\textbf{Oracle} (\texttt{expose\_\allowbreak{}realised}) & Revealing the realised delay to the agent (an information upper bound). \\
\textbf{U\{0..4\}} & Delay drawn uniformly from \{0, 1, 2, 3, 4\} (exact discrete uniform, \texttt{uniform\_int}, amendment A-4 of Phase 3). \\
\textbf{C0} & No delay. \\
\textbf{C1} & Fixed d = 1 + queue. The standard condition. (Phase 3 Stage 1 called it \texttt{fixed1-aug}; Phase 4 B/C/T use it without previous action.) \\
\textbf{C2} & Fixed d = 1, no queue (Phase 2's augmentation ablation). \\
\textbf{C3} & d \textasciitilde{} U\{0..4\} redrawn every step, queue shown. Turned out to leave nothing to infer (A-1). \\
\textbf{C4} & d \textasciitilde{} U\{0..4\} redrawn every step, queue hidden. Memory has something to recover (G0 memory headroom). \\
\textbf{C5} & Time-varying (non-stationary) delay. Optional in the Phase 3 plan; not run. \\
\textbf{C6} & d \textasciitilde{} U\{0..4\} drawn once per episode and held, queue shown (the G2s condition). \textbf{C6o}: the same with the oracle. \\
\textbf{C1pa / C4pa} & Phase 4 Stage A labels for C1 / C4 with the previous action. \\
\textbf{Headroom arms A1--A4} (Phase 3 S0.1) & A1 = fixed d = 1 + queue; A2 = stochastic + queue; A3 = A2 + oracle; A4 = stochastic, no queue. \\
\textbf{Sufficient statistic / Markov} & An observation that contains everything needed to predict the future. Delay breaks it; the queue restores it. \\
\bottomrule
\end{tabular}\end{center}



\section*{6. Confounds}\addcontentsline{toc}{section}{6. Confounds}



\begin{center}\small\begin{tabular}{>{\raggedright\arraybackslash}p{0.128\linewidth}>{\raggedright\arraybackslash}p{0.505\linewidth}>{\raggedright\arraybackslash}p{0.274\linewidth}}
\toprule
Label & Meaning & Removed by \\
\midrule
\textbf{C-1, capacity} & Every recursive trunk had \textasciitilde{}788k parameters vs the MLP's 135k. & Phase 3 Stage 1 (G1) \\
\textbf{C-2, latent reset} & The port re-initialised its latent every call (HRM-Agent's weaker "Reset-Z" variant). & Phase 3 Stage 2 Carry-Z (G3) \\
\textbf{C-3, nothing to infer} & Under fixed delay + queue the observation is already sufficient. & Phase 3 G0 headroom, Stage 2 C4, C6 (G2, G2s) \\
\bottomrule
\end{tabular}\end{center}



\section*{7. Trunks (arms)}\addcontentsline{toc}{section}{7. Trunks (arms)}



Parameter counts are actor + critic under DDPG on Pendulum with the queue, unless stated.



\begin{center}\small\begin{tabular}{>{\raggedright\arraybackslash}p{0.237\linewidth}>{\raggedright\arraybackslash}p{0.586\linewidth}>{\raggedleft\arraybackslash}p{0.083\linewidth}}
\toprule
Config & What it is & Params \\
\midrule
\textbf{\texttt{dnn}} (MLP 135k) & MLP, hidden [256, 256], ReLU & 134,914 \\
\textbf{\texttt{dnn\_large}} (MLP 788k) & MLP, hidden [624, 624] & 788,114 \\
\textbf{\texttt{trm}} (TRM 788k) & TRM, hidden 128, H = 3, L = 6, 2-layer SwiGLU cell & 787,842 \\
\textbf{\texttt{trm\_small}} (TRM 136k) & TRM, hidden 44 & 135,654 \\
\textbf{\texttt{trm\_h1}} & \texttt{trm} with H = 1 (7 cell applications instead of 21) & 787,842 \\
\textbf{\texttt{trm\_h6}} & \texttt{trm} with H = 6 (42 applications) & 787,842 \\
\textbf{\texttt{trm\_hier6}} & Phase 2 variant: six parallel low-level latents & \textasciitilde{}794k \\
\textbf{\texttt{trm\_singlez}} & Phase 2 variant: a single latent, no z\_H / z\_L split & --- \\
\textbf{\texttt{trm\_carryz}} (Carry-Z) & \texttt{trm} that carries its latent across steps (stored-state replay) & = \texttt{trm} \\
\textbf{Reset-Z} & The ordinary \texttt{trm}: latent re-initialised every call & = \texttt{trm} \\
\textbf{\texttt{trm\_fullgrad}} & \texttt{trm} + \texttt{full\_grad}: gradient through all 21 applications & = \texttt{trm} \\
\textbf{\texttt{trm\_\allowbreak{}fullgrad\_\allowbreak{}ds}} & \texttt{trm\_fullgrad} + \texttt{deep\_sup}: critic supervised on every outer cycle & = \texttt{trm} \\
\textbf{\texttt{transformer}} & Self-attention trunk (Phases 1--2) & \textasciitilde{}797k \\
\textbf{\texttt{lstm}} & LSTM, stateless across calls (Phases 1--2) & \textasciitilde{}167k \\
\textbf{\texttt{lstm\_carry}} & LSTM with its hidden state carried (Phase 3, untuned) & --- \\
\textbf{TRM-trunc} & Phase 4 name for truncated \texttt{trm} with clipping (Stage A baseline) & = \texttt{trm} \\
\bottomrule
\end{tabular}\end{center}



Under TD3 (two critics) the "788k" configs total about 1.18M and the "135k" configs about 202k.



\section*{8. TRM internals}\addcontentsline{toc}{section}{8. TRM internals}



\begin{center}\small\begin{tabular}{>{\raggedright\arraybackslash}p{0.179\linewidth}>{\raggedright\arraybackslash}p{0.755\linewidth}}
\toprule
Term & Meaning \\
\midrule
\textbf{TRM} & Tiny Recursive Model (Jolicoeur-Martineau, 2025): a small shared cell applied many times instead of adding parameters. \\
\textbf{Cell / reasoning module} & The shared network applied every step: 2 SwiGLU blocks with RMSNorm residuals. Attention was dropped in the port (a flat observation is a length-1 sequence). \\
\textbf{z\_H, z\_L} & The high-level ("answer") and low-level ("reasoning") latents. \\
\textbf{H\_cycles (H)} & Outer cycles. Each updates z\_H once. \\
\textbf{L\_cycles (L)} & Inner cycles per outer cycle. Each updates z\_L. \\
\textbf{21 applications} & H $\times$ (L + 1) = 3 $\times$ 7 cell applications per forward pass at H = 3. \\
\textbf{One-step gradient truncation} & All outer cycles but the last run under \texttt{no\_grad}; \textbf{only 7 of 21 applications receive gradient}. \\
\textbf{\texttt{full\_grad}} & Phase 4 flag removing the truncation (all 21 receive gradient). \\
\textbf{\texttt{deep\_sup}} / deep supervision & Phase 4 flag: the critic's TD loss is applied to every outer cycle's readout (Lee et al., 2015). TRM's original deep supervision re-feeds the carry over many steps against a fixed answer. \\
\textbf{ACT / halting} & Adaptive computation time: TRM learns when to stop iterating. Dropped in the port (no analogue in off-policy RL). \\
\textbf{H\_init, L\_init} & The fixed initial latents. \\
\textbf{Carry} & The latent passed between decision steps in Carry-Z. \\
\textbf{SwiGLU / RMSNorm} & The gated MLP block and the normalisation TRM uses. \\
\textbf{embed\_scale} & The input projection is multiplied by $\surd$hidden\_size (needed for the C export). \\
\bottomrule
\end{tabular}\end{center}



\section*{9. Reinforcement-learning terms}\addcontentsline{toc}{section}{9. Reinforcement-learning terms}



\begin{center}\small\begin{tabular}{>{\raggedright\arraybackslash}p{0.206\linewidth}>{\raggedright\arraybackslash}p{0.728\linewidth}}
\toprule
Term & Meaning \\
\midrule
\textbf{Model-free RL} & Learning a policy and value function from sampled transitions, never using the plant's equations. \\
\textbf{Actor / critic} & The policy network (observation $\rightarrow$ action) and the value network (observation, action $\rightarrow$ Q). \\
\textbf{DDPG} & Deep Deterministic Policy Gradient (Lillicrap et al., 2016). Phases 1--3 and Phase 4 Stages A and C. Actor LR 1e-4, critic LR 1e-3, batch 128, warmup 1000. \\
\textbf{TD3} & Twin Delayed DDPG (Fujimoto et al., 2018): twin critics, delayed actor updates, target smoothing. Phase 4 Stages B and T. \texttt{td3.yaml}: LR 1e-3 / 1e-3, batch 256, warmup 10,000. \\
\textbf{SAC, PPO} & Other algorithms in the codebase; never used for results. \\
\textbf{OU noise} & Ornstein--Uhlenbeck exploration noise used by DDPG. \\
\textbf{Replay buffer} & Stored transitions sampled for off-policy updates. \\
\textbf{Warmup} & Initial steps with random actions. \\
\textbf{TD target, critic loss} & The bootstrapped value target and the critic's squared error against it. \\
\textbf{Value fit} & How well the critic fits; measured as late critic loss (median over the last quarter of episodes). \\
\textbf{Stored-state replay (R2D2-style)} & Saving the carried latent with each transition so updates start from the state the policy acted from (Kapturowski et al., 2019). Verified exact for Carry-Z. \\
\textbf{Burn-in} & Alternative to stored state (recompute the carry over a prefix). Not used. \\
\textbf{Gradient clipping} & Capping the gradient norm (Phase 4 Stage A: actor 15, critic 2800, from the a0 pilot's p99). \\
\textbf{Critic divergence} & Critic loss rising instead of falling over training (TRM 788k under TD3 in Stage T). \\
\bottomrule
\end{tabular}\end{center}



\section*{10. Environments}\addcontentsline{toc}{section}{10. Environments}



\begin{center}\small\begin{tabular}{>{\raggedright\arraybackslash}p{0.287\linewidth}>{\raggedright\arraybackslash}p{0.072\linewidth}>{\raggedright\arraybackslash}p{0.106\linewidth}>{\raggedright\arraybackslash}p{0.413\linewidth}}
\toprule
Environment & Obs / act dims & "Solved" reference & Status in this project \\
\midrule
\textbf{Pendulum-v1} & 3 / 1 & $\approx$ $-$150 & The only task every trunk learned. All decisive tests ran here. \\
\textbf{MountainCarContinuous-v0} & 2 / 1 & +90 & Floor effect (Phase 1): exploration-limited under DDPG. \\
\textbf{BipedalWalker-v3} & 24 / 4 & +300 & Floor effect (Phase 2): no run learned to walk. \\
\textbf{InvertedDoublePendulum-v5 (IDP)} & 9 / 1 & 9100 (bar 8190) & Phase 4: MLP not reliable with delay (GB0 fail); learns without delay but doesn't hold it. \\
\textbf{Hopper-v5} & 11 / 3 & 3800 (bar 3420) & Phase 4 fallback: fail; dropped. \\
\textbf{Floor effect} & --- & --- & Nothing learns, so every trunk scores the same and comparisons mean nothing. \\
\textbf{Complexity rung} & --- & --- & A harder task, used to test H1. \\
\bottomrule
\end{tabular}\end{center}



\section*{11. Metrics and statistics}\addcontentsline{toc}{section}{11. Metrics and statistics}



\begin{center}\small\begin{tabular}{>{\raggedright\arraybackslash}p{0.217\linewidth}>{\raggedright\arraybackslash}p{0.717\linewidth}}
\toprule
Term & Meaning \\
\midrule
\textbf{avg-100 / final return} & Mean training return over the last 100 episodes; the primary metric throughout. Includes exploration noise. \\
\textbf{eval\_return / eval last-5} & A noise-free evaluation episode every 10 training episodes; the mean of the last 5. Secondary. \\
\textbf{Battery} & Phase 3's evaluation harness (\texttt{w0\_eval.py}): 20 deterministic episodes on unseen seeds plus grids for delay mismatch, sensor noise, plant mass/length, actuator dropout; control metrics. \\
\textbf{IQM} & Interquartile mean (Agarwal et al., 2021). Used for battery returns. \\
\textbf{Stratified bootstrap CI} & Confidence interval resampling runs then episodes. \\
\textbf{Sign test} & Counts how many paired differences are positive; no distributional assumption. \\
\textbf{k/n} & k of n paired differences positive (e.g. 0/12). \\
\textbf{One-sided / two-sided p} & 5/6 $\rightarrow$ one-sided 0.109; 6/6 $\rightarrow$ 0.016; 12/12 $\rightarrow$ two-sided 0.0005; 10/12 $\rightarrow$ two-sided 0.039. \\
\textbf{Pair} & Two runs matched by seed (and condition, when pooled to n = 12). \\
\textbf{Pooled SD} & $\surd$((var\_a + var\_b)/2) across seeds; the magnitude threshold in G0, GC, G2s. \\
\textbf{Span, untrained\_E, best\_E} & best $-$ untrained return on an environment; normalises deficits for H1$'$. Untrained = uniform-random policy (R-2: Pendulum $-$1164, IDP 49.9, Hopper 16.3). Phases 1--3 quoted $\approx$ $-$1460 for an untrained Pendulum \emph{actor}. \\
\textbf{$\delta$\_E(s)} & Normalised deficit: (MLP $-$ TRM) / span on environment E for seed s. \\
\textbf{"Learned"} & B0's bar: avg-100 $\geq$ 90\% of the reference return. \\
\textbf{"Clear of floor"} & Draft K-3's exploratory bar ($\geq$ 2275 on IDP). Superseded. \\
\textbf{Reversals/s} & Torque direction changes per second in deterministic rollouts; smoothness. \\
\textbf{abs($\Delta$u), effort, saturation, hf\_frac, work} & Command rate, control energy, fraction at the limit, high-frequency power share, mechanical work. \\
\textbf{Noise drop} & Nominal IQM $-$ IQM under sensor noise $\sigma$ = 0.05 of the observation half-range. \\
\textbf{Lipschitz constant} & Spectral norm of $\partial$a/$\partial$s, how sharply the actor reacts to its input (Stage E). \\
\textbf{Latency p50 / p99} & Median / 99th-percentile time per decision (batch 1, 1 CPU thread). \\
\textbf{Matched episodes / wall-clock / latency} & The three ways Phase 3 reported cost. \\
\textbf{Clip fraction} & Share of updates where gradient clipping fired. \\
\bottomrule
\end{tabular}\end{center}



\section*{12. Amendments}\addcontentsline{toc}{section}{12. Amendments}



\subsection*{Phase 3 (recorded 2026-09-21, before data)}\addcontentsline{toc}{subsection}{Phase 3 (recorded 2026-09-21, before data)}



\begin{center}\small\begin{tabular}{>{\raggedright\arraybackslash}p{0.092\linewidth}>{\raggedright\arraybackslash}p{0.842\linewidth}}
\toprule
Label & Change \\
\midrule
\textbf{A-1} & Theory check: with per-step redrawn delay and the queue, the observation is already Markov, so the oracle arm must be null. \\
\textbf{A-2} & Adds arm A4 and lets G0 pass on memory headroom (A2 $-$ A4). \\
\textbf{A-3} & C4 becomes co-primary with C3 for G2. \\
\textbf{A-4} & Exact discrete-uniform delay sampling. \\
\textbf{A-5} & Every job on a GPU node; battery as 12 shard jobs. \\
\textbf{A-6} & Every Stage 2 trunk also sees its previous action. \\
\bottomrule
\end{tabular}\end{center}



\subsection*{Phase 4 (each dated, before the data it affects)}\addcontentsline{toc}{subsection}{Phase 4 (each dated, before the data it affects)}



\begin{center}\small\begin{tabular}{>{\raggedright\arraybackslash}p{0.087\linewidth}>{\raggedright\arraybackslash}p{0.087\linewidth}>{\raggedright\arraybackslash}p{0.733\linewidth}}
\toprule
Label & Date & Change \\
\midrule
\textbf{R-1} & 10-06 & Stage A reuses Phase 3's MLP runs (bit-identity check passed). \\
\textbf{R-2} & 10-06 & Random-policy baselines frozen for the H1$'$ span. \\
\textbf{V-1} & 10-07 & Code verification passed (30/30 bit-identical; gradient-hook test). \\
\textbf{A-1} & 10-07 & Stage A clip norms: actor 15, critic 2800. \\
\textbf{K-1} & 10-07 & Hopper kill switch fired. \\
\textbf{K-2} & 10-07 & Hopper dropped; IDP diagnostic added. \\
\textbf{T-1} & 10-07 & Stage T added (both conditions, \texttt{td3.yaml} as-is). \\
\textbf{K-3} & 10-08 & B1 not run; H1$'$ not run; exploratory draft superseded. \\
\bottomrule
\end{tabular}\end{center}



\section*{13. Other code lists you may meet}\addcontentsline{toc}{section}{13. Other code lists you may meet}



\begin{center}\small\begin{tabular}{>{\raggedright\arraybackslash}p{0.212\linewidth}>{\raggedright\arraybackslash}p{0.088\linewidth}>{\raggedright\arraybackslash}p{0.605\linewidth}}
\toprule
Codes & Where & Meaning \\
\midrule
\textbf{E1--E4} & Phase 1--2 report \S{}7 & Mechanistic explanations: E1 optimisation/conditioning deficit (supported), E2 slower learning, E3 lower asymptote, E4 instability. \\
\textbf{E1--E5} & Phase 1--2 report \S{}10 & Recommended experiments: E1 capacity-matched MLP, E2 Carry-Z, E3 depth sweep, E4 longer budgets, E5 restore per-cycle supervision. They became Phase 3 Stage 1/2 and Phase 4 Stage A. \\
\textbf{E-1--E-9} & Phase 3 plan \S{}7 & Engineering items (evaluation harness, Carry-Z, perturbation wrappers, oracle arm, configs, latency benchmark, MuJoCo, deep-supervision variants, observation-delay wrappers). \\
\textbf{A1--A9, B1--B10, C1--C8, D1--D5, E1--E4, F1--F5, G1--G4} & Phase 3 plan \S{}6 & The exploration inventory (mechanism, conditioning, control practicality, tasks, learner, methodology, out of scope). Not used elsewhere. \\
\bottomrule
\end{tabular}\end{center}



\section*{14. Infrastructure}\addcontentsline{toc}{section}{14. Infrastructure}



\begin{center}\small\begin{tabular}{>{\raggedright\arraybackslash}p{0.351\linewidth}>{\raggedright\arraybackslash}p{0.583\linewidth}}
\toprule
Term & Meaning \\
\midrule
\textbf{Katana} & UNSW's HPC cluster (PBS Pro scheduler), where every Phase 1--4 job ran. \\
\textbf{PBS, qsub, qstat, qdel} & The scheduler and its submit / status / delete commands. \\
\textbf{mechgpu2 / mechgpu12} & GPU queues; $\leq$ 2 h jobs reach every GPU node, longer ones only group-owned nodes. \\
\textbf{Walltime} & A job's time limit (1:59 for most Phase 3--4 jobs). \\
\textbf{Frozen snapshot} (\texttt{p3code/\allowbreak{}<stage>}, \texttt{p4code/\allowbreak{}<stage>}) & A copy of the code per stage; \texttt{sync.sh} refuses to overwrite one jobs were submitted from. \\
\textbf{Ledger} & \texttt{\textasciitilde{}/\allowbreak{}jobs/\allowbreak{}phase\{3,\allowbreak{}4\}/\allowbreak{}ledger.\allowbreak{}jsonl}: every submitted job. \\
\textbf{\texttt{p3.py} / \texttt{p4.py}} & The job managers (manifest, submit, status, resubmit). \\
\textbf{\texttt{gates.json} / \texttt{gates\_eval.py}} & The frozen rules and the scripts that evaluate them. \\
\textbf{Run board} & The live status artifact (Phase 3 and Phase 4 each have one). \\
\textbf{model\_id} & A unique ID per run; the checkpoint is \texttt{models/\allowbreak{}<model\_\allowbreak{}id>.\allowbreak{}pt}. \\
\textbf{Exit $-$19} & A PBS code reported after two Phase 4 jobs had finished cleanly; counted done. \\
\textbf{k207 / RTX PRO 6000} & A Blackwell GPU node with no kernels in the venv's PyTorch build; excluded with \texttt{gpu\_\allowbreak{}rp6000=False}. \\
\textbf{Bit-identical} & Same weights and histories to the last bit; the standard for "unchanged code path". \\
\bottomrule
\end{tabular}\end{center}



\section*{15. Literature you will see cited}\addcontentsline{toc}{section}{15. Literature you will see cited}



\begin{center}\small\begin{tabular}{>{\raggedright\arraybackslash}p{0.149\linewidth}>{\raggedright\arraybackslash}p{0.293\linewidth}>{\raggedright\arraybackslash}p{0.463\linewidth}}
\toprule
Short name & Work & Why it matters here \\
\midrule
\textbf{TRM} & Jolicoeur-Martineau (2025), \emph{Less is More: Recursive Reasoning with Tiny Networks} & The architecture under test. \\
\textbf{HRM} & Hierarchical Reasoning Model, TRM's predecessor & TRM simplifies its two networks into one cell. \\
\textbf{HRM-Agent} & Dang \& Rawlinson (2025) & Closest precedent: RL-trained recursive model on discrete mazes; source of Carry-Z vs Reset-Z. \\
\textbf{TRC} & Jain \& Linares (2025), \emph{Tiny Recursive Control} & Model-based TRM for control; each pass gets dynamics-error feedback. The contrast that explains why the port's depth goes unused. \\
\textbf{TRM-PL} & Wijayatunga \& Armellin (2026) & TRM for the J2-perturbed Lambert problem; also model-based. \\
\textbf{ARC Prize (2025)} & \emph{The Hidden Drivers of HRM's Performance} & Training-time refinement, not inference depth, drives HRM. \\
\textbf{Roye-Azar et al. (2026)} & TRM on ARC-AGI-1 & Most accuracy arrives at the first recursion step; removing the puzzle ID collapses accuracy. \\
\textbf{Katsikopoulos \& Engelbrecht (2003)} & Delayed MDPs & The queue-augmentation theorem. \\
\textbf{Derman et al. (2021)} & Non-stationary delays & Why time-varying delay (C5) is the hardest case. \\
\textbf{Ni et al. (2022)} & Recurrent model-free RL for POMDPs & Previous-action input (A-6); tuned recurrent baselines. \\
\textbf{R2D2} & Kapturowski et al. (2019) & Stored-state replay for carried latents. \\
\textbf{DSN} & Lee et al. (2015), Deeply-Supervised Nets & The deep-supervision surrogate in Stage A. \\
\textbf{DEQ} & Bai et al. (2019), Deep Equilibrium Models & Alternative to truncation; not run. \\
\textbf{Agarwal et al. (2021)} & Statistical precipice & IQM and stratified bootstrap. \\
\textbf{Matheron et al. (2019)} & The problem with DDPG & Why MountainCar was a floor. \\
\textbf{Fujimoto et al. (2018)} & TD3 & The second algorithm; IDP reference curves. \\
\textbf{Lillicrap et al. (2016)} & DDPG & The main algorithm. \\
\bottomrule
\end{tabular}\end{center}
\end{document}
